\section{Gaussian decoding details}
\label{sec:supp-codec}

A Gaussian stores its center $\boldsymbol\mu$, covariance $\boldsymbol\Sigma$, opacity $\alpha$, and directional appearance.

Background, rigid, and human Gaussian decoders independently map each query to $K$ child Gaussians. For child $k$ of query $i$, the center is confined to its support cell:

\begin{equation}
\boldsymbol\mu_{i,k}=\mathbf q_i+
\tfrac12\boldsymbol\delta_i\odot\tanh(\mathbf b_{i,k}).
\label{eq:child-center}
\end{equation}

Positive scales and normalized rotations determine positive-definite covariance. Opacity is multiplied by $1-f_i$, so free queries cannot produce visible Gaussians. Human children inherit the skinning weights of their support query; rigid children share their asset's rigid transform. A separate sky head predicts cubemap colors, with the observation-supported domain retained in $V_{\rm sky}$.

Appearance uses a constant RGB term and three local directional coefficients per channel, evaluated as $c(d)=\max(0,c_0+C_{\rm app}\hat d_{\rm local})$. Decoding is per query: once the codec is frozen, an unchanged observed code produces unchanged Gaussian parameters. The completion network can exchange context across queries without changing this decoding property.
